RESTORE treats its address-role EA as the base of the save area without first reading an EA-sized operand. It restores R0--R15 and FLAGS from the fixed base block. A set \texttt{GS\_VALID} bit selects the corresponding GS image for validation and restoration; a clear bit preserves the current GSn. User-mode RESTORE ignores the saved STATUS image, while supervisor RESTORE applies the normal STATUS write rules. A set state-block bitmap bit restores its CPUID-described extension component. A clear bit for a described component installs that component's initial state and makes the component clean; a restored non-initial image makes the component modified.

Before reading any extension component or changing architectural state, RESTORE validates the complete base header and bitmap. An unrecognized format, a nonzero reserved bit, or any set bitmap bit without a matching CPUID component descriptor raises INVALID\_CONTROL\_STATE. RESTORE does not ignore, skip, or zero-fill unknown components.

The complete \texttt{SAVE\_AREA\_SIZE\_BYTES} range reported by CPUID, every selected GS image, and every selected extension component are validated before architectural state is committed. Any address, access, or validation fault leaves the architectural register state and save area unchanged.
